\documentclass[12pt]{article}
\usepackage{amsmath, amsthm, amssymb}
\usepackage{geometry}
\usepackage{hyperref}
\usepackage{booktabs}
\geometry{margin=1in}

\newtheorem{theorem}{Theorem}
\newtheorem{lemma}[theorem]{Lemma}
\newtheorem{corollary}[theorem]{Corollary}
\newtheorem{observation}[theorem]{Observation}
\newtheorem{conjecture}[theorem]{Conjecture}
\newtheorem{remark}[theorem]{Remark}
\theoremstyle{definition}
\newtheorem{definition}[theorem]{Definition}

\title{The Staircase Eigenspace Property:\\
Computational Verification and Structural Analysis}
\author{Clio}
\date{2026-04-15}

\begin{document}
\maketitle

\begin{abstract}
Let $\Pi = \prod_j (1 + s_{i_j})$ be the staircase symmetrizing product in the group algebra of $S_n$, and let $H = \langle s_1, s_3, s_5, \ldots \rangle$ with $V^H$ the joint $+1$ eigenspace (fixed space) of all generators of $H$ in the regular representation. We test the claim that after applying the staircase cascade through blocks $1, \ldots, j$, the resulting image $\Phi_j(V^H)$ lies in the $+1$ eigenspace of $s_j$ (i.e., $s_j$ fixes every vector in the image). Using exact rational arithmetic in the regular representation of $S_n$, we verify this for $n = 4, 5, 6, 7$. Strikingly, the claim holds not just for odd $j$ (as originally conjectured) but for \emph{every} $j$, both odd and even. Moreover, we discover a refined structure: at step $j$, additional generators $s_{n-1}, s_{n-2}, \ldots, s_{j+1}$ (those not yet ``consumed'' by the cascade) also fix the image. The image dimension equals $\dim(V^H) = n!/2^{\lfloor n/2 \rfloor}$ throughout---no collapsing occurs.
\end{abstract}

\section{Setup and Definitions}

\subsection{Symmetric group and generators}
Let $S_n$ denote the symmetric group on $n$ letters. We represent permutations in one-line notation as tuples $(\sigma(1), \ldots, \sigma(n))$. The adjacent transposition $s_i$ ($1 \le i \le n-1$) acts by \emph{left multiplication} on one-line notation: swapping the values in positions $i$ and $i+1$.

More precisely, if $\sigma = (\sigma(1), \ldots, \sigma(n))$, then
\[
  s_i \cdot \sigma = (\sigma(1), \ldots, \sigma(i-1), \sigma(i+1), \sigma(i), \sigma(i+2), \ldots, \sigma(n)).
\]

The group algebra $\mathbb{Q}[S_n]$ has basis $\{e_\sigma : \sigma \in S_n\}$, and $s_i$ acts by $s_i \cdot e_\sigma = e_{s_i \sigma}$.

\subsection{The subgroup $H$ and fixed space $V^H$}
Define
\[
  H = \langle s_1, s_3, s_5, \ldots \rangle \cong S_{\lfloor n/2 \rfloor} \times S_{\lceil n/2 \rceil - \lfloor n/2 \rfloor} \cdots
\]
(the subgroup generated by the odd-indexed adjacent transpositions). In the regular representation $\mathbb{Q}[S_n]$, the \emph{fixed space} is:
\[
  V^H = \{v \in \mathbb{Q}[S_n] : s_i \cdot v = v \text{ for all odd } i\}.
\]

\begin{observation}[Dimension of $V^H$]
  \[
    \dim(V^H) = \frac{n!}{2^{\lfloor n/2 \rfloor}}.
  \]
  This equals $|S_n|/|H|$ and is confirmed computationally.
\end{observation}

\begin{proof}
  $H \cong (\mathbb{Z}/2\mathbb{Z})^{\lfloor n/2 \rfloor}$, so $|H| = 2^{\lfloor n/2 \rfloor}$.
  The fixed space $V^H$ in the regular representation is the space of $H$-averages, which has
  $\dim(V^H) = n!/|H| = n!/2^{\lfloor n/2 \rfloor}$, as $H$ acts freely (it embeds in $S_n$ as a
  direct product of disjoint transpositions, which acts freely on $S_n$ by left multiplication).
\end{proof}

\subsection{The staircase cascade}
Define:
\begin{itemize}
  \item $\text{block}_j = (1+s_j)(1+s_{j-1})\cdots(1+s_1) \in \mathbb{Q}[S_n]$
  \item $\Phi_j = \text{block}_1 \cdot \text{block}_2 \cdots \text{block}_j$ (left-multiplication composition)
\end{itemize}
Since $(1+s_i)^2 = 2(1+s_i)$ (as $s_i^2 = 1$), each factor $(1+s_i)$ appears idempotently in repeated applications. The cascade applied to a vector $v$ is computed iteratively:
\[
  v \xrightarrow{\text{block}_1} (1+s_1)v \xrightarrow{\text{block}_2} (1+s_2)(1+s_1)\bigl[(1+s_1)v\bigr] \xrightarrow{\cdots}
\]
Up to positive scalar factors of $2^{j(j-1)/2}$, we have:
\[
  \Phi_j(v) \sim (1+s_j)(1+s_{j-1})\cdots(1+s_1) \cdot v.
\]

\section{The Main Claim}

\begin{conjecture}[Staircase Eigenspace Property]
  For all $j = 1, \ldots, n-1$ and all $v \in V^H$:
  \[
    s_j \cdot \Phi_j(v) = \Phi_j(v).
  \]
  That is, $\Phi_j(V^H) \subseteq \ker(s_j - \mathrm{id})$.
\end{conjecture}

The original statement asked this only for \emph{odd} $j$, placing the image in the $+1$ eigenspace of odd generators. The computation below shows it holds for \emph{all} $j$.

\section{Computational Verification}

\subsection{Method}

We work entirely in the regular representation of $S_n$ over $\mathbb{Q}$, using exact \texttt{fractions.Fraction} arithmetic throughout. Concretely:
\begin{enumerate}
  \item Enumerate all $n!$ permutations in lexicographic order.
  \item Represent subspaces as lists of sparse vectors (dictionaries $\sigma \mapsto c_\sigma \in \mathbb{Q}$).
  \item Compute $V^H$ by sequentially applying the projector $(1+s_i)/2$ for each odd-indexed generator $s_i$ of $H$.
  \item Apply each staircase block $(1+s_j)$ iteratively, reducing to a basis after each step using Gaussian elimination.
  \item After block $j$, check: does $s_j \cdot w = w$ for every basis vector $w$ in the image?
\end{enumerate}

The Gaussian elimination is exact (no floating-point), and all intermediate values are rational numbers.

\subsection{Results}

\begin{table}[h]
\centering
\caption{Staircase cascade results for $n=4$. All s\_j fix the image (YES).}
\begin{tabular}{ccccc}
\toprule
Block $j$ & Parity & $\dim(\Phi_j(V^H))$ & $s_j$ fixes image & $\Phi_j(V^H) = V^H$? \\
\midrule
1 & odd  & 6 & YES & YES \\
2 & even & 6 & YES & NO \\
3 & odd  & 6 & YES & NO \\
\bottomrule
\end{tabular}
\end{table}

\begin{table}[h]
\centering
\caption{Staircase cascade results for $n=5$.}
\begin{tabular}{ccccc}
\toprule
Block $j$ & Parity & $\dim(\Phi_j(V^H))$ & $s_j$ fixes image & $\Phi_j(V^H) = V^H$? \\
\midrule
1 & odd  & 30 & YES & YES \\
2 & even & 30 & YES & NO \\
3 & odd  & 30 & YES & NO \\
4 & even & 30 & YES & NO \\
\bottomrule
\end{tabular}
\end{table}

\begin{table}[h]
\centering
\caption{Staircase cascade results for $n=6$.}
\begin{tabular}{ccccc}
\toprule
Block $j$ & Parity & $\dim(\Phi_j(V^H))$ & $s_j$ fixes image & $\Phi_j(V^H) = V^H$? \\
\midrule
1 & odd  & 90 & YES & YES \\
2 & even & 90 & YES & NO \\
3 & odd  & 90 & YES & NO \\
4 & even & 90 & YES & NO \\
5 & odd  & 90 & YES & NO \\
\bottomrule
\end{tabular}
\end{table}

\begin{table}[h]
\centering
\caption{Staircase cascade results for $n=7$.}
\begin{tabular}{ccccc}
\toprule
Block $j$ & Parity & $\dim(\Phi_j(V^H))$ & $s_j$ fixes image & $\Phi_j(V^H) = V^H$? \\
\midrule
1 & odd  & 630 & YES & YES \\
2 & even & 630 & YES & NO \\
3 & odd  & 630 & YES & NO \\
4 & even & 630 & YES & NO \\
5 & odd  & 630 & YES & NO \\
6 & even & 630 & YES & NO \\
\bottomrule
\end{tabular}
\end{table}

\begin{observation}[Universal eigenspace property]
  For $n = 4, 5, 6, 7$, the staircase eigenspace property holds for \textbf{every} $j$, both odd and even:
  \[
    s_j \cdot \Phi_j(V^H) = \Phi_j(V^H) \quad \text{for all } j = 1, \ldots, n-1.
  \]
\end{observation}

\begin{observation}[Dimension preservation]
  $\dim(\Phi_j(V^H)) = \dim(V^H) = n!/2^{\lfloor n/2 \rfloor}$ for all $j$.
  The staircase operators do not collapse the image.
\end{observation}

\subsection{Refined pattern: which other generators fix the image?}

Beyond the main claim, we tracked which \emph{other} generators $s_k$ (for $k \ne j$) fix the image $\Phi_j(V^H)$ at each step. The data reveals a striking tail structure.

\begin{table}[h]
\centering
\caption{$n=6$: which $s_k$ fix $\Phi_j(V^H)$? YES = fixes, no = does not fix, (*) = $k=j$ (main claim).}
\begin{tabular}{c|cccccc}
\toprule
$j$ & $s_1$ & $s_2$ & $s_3$ & $s_4$ & $s_5$ \\
\midrule
initial $V^H$ & YES & no & YES & no & YES \\
\midrule
1 & (*) & no & YES & no & YES \\
2 & no  & (*) & no & no & YES \\
3 & no  & no  & (*) & no & YES \\
4 & no  & no  & no  & (*) & no \\
5 & no  & no  & no  & no  & (*) \\
\bottomrule
\end{tabular}
\end{table}

\begin{table}[h]
\centering
\caption{$n=7$: which $s_k$ fix $\Phi_j(V^H)$?}
\begin{tabular}{c|cccccc}
\toprule
$j$ & $s_1$ & $s_2$ & $s_3$ & $s_4$ & $s_5$ \\
\midrule
initial $V^H$ & YES & no & YES & no & YES \\
\midrule
1 & (*) & no & YES & no & YES \\
2 & no  & (*) & no & no & YES \\
3 & no  & no  & (*) & no & YES \\
4 & no  & no  & no  & (*) & no \\
5 & no  & no  & no  & no  & (*) \\
6 & no  & no  & no  & no  & no \\
\bottomrule
\end{tabular}
\end{table}

\begin{observation}[Tail symmetry pattern]
  Let $n^* = n-1$ (the index of the last generator). In the data for $n = 6, 7$:
  \begin{itemize}
    \item At step $j = 1$: generators $s_3, s_5$ (the odd generators $\ge 3$) still fix the image.
    \item At step $j = 2$: $s_5$ fixes the image.
    \item At step $j = 3$: $s_5$ still fixes the image.
    \item At step $j = 4$: only $s_j = s_4$ fixes (among generators $\le 5$).
    \item At step $j = 5$ ($n=6$): only $s_5$.
  \end{itemize}

  For $n=6$: $s_{n-1} = s_5$ fixes the image at steps $j = 1, 2, 3$, then stops at $j = 4$.
  Pattern: $s_{n-1}$ fixes $\Phi_j(V^H)$ for $j \le n-2$, not for $j = n-1$.

  More precisely (for $n=6, 7$): the set of generators fixing $\Phi_j(V^H)$ includes $s_j$ (the main claim) and the ``tail'' generators $s_k$ for $k > j$ that were generators of $H$ in $V^H$ AND have not yet been disturbed by the cascade.
\end{observation}

\section{Theoretical Explanation}

\subsection{Why $(1+s_j)$ creates a $+1$ eigenvector for $s_j$}

The key algebraic identity: for any $s_j$ with $s_j^2 = 1$,
\[
  s_j \cdot (1 + s_j) = s_j + s_j^2 = s_j + 1 = (1 + s_j).
\]
Therefore, $(1+s_j)w$ is a $+1$ eigenvector of $s_j$ for \emph{any} $w$.

This immediately explains why the image of applying $(1+s_j)$ to anything lies in the $+1$ eigenspace of $s_j$. But our cascade does \emph{not} simply apply $(1+s_j)$ as the last step. Block $j$ applies $(1+s_j)(1+s_{j-1})\cdots(1+s_1)$, but then \emph{block $j+1$} later applies $(1+s_{j+1})\cdots(1+s_1)$ again---including a fresh $(1+s_1)$ which could disturb the $s_j$ eigenspace.

\subsection{Why the image stays in the $+1$ eigenspace of $s_j$ after subsequent blocks}

After block $j$, subsequent blocks $j+1, j+2, \ldots$ apply operators that include $(1+s_k)$ for $k < j$. The question is: does applying $(1+s_{j-1}), (1+s_{j-2}), \ldots$ to a vector in the $+1$ eigenspace of $s_j$ keep it there?

The answer depends on the commutation relations in $S_n$:
\begin{itemize}
  \item $s_j s_k = s_k s_j$ when $|j-k| \ge 2$ (far commutativity).
  \item $s_j s_{j\pm 1} s_j = s_{j\pm 1} s_j s_{j\pm 1}$ (braid relation).
\end{itemize}

For $|j - k| \ge 2$: if $s_j w = w$, then $s_j (1+s_k)w = (1+s_k) s_j w = (1+s_k)w$, so the $+1$ eigenspace of $s_j$ is preserved by $(1+s_k)$. This covers many cases.

The braid relation case ($k = j \pm 1$) is more subtle. The computation shows empirically that even these cases preserve the eigenspace property---but explaining this requires understanding the structure of $\Phi_j(V^H)$ more carefully.

\subsection{Connection to the rank hierarchy}

The space $V^H$ corresponds to the image of the product $\Pi_{H} = \prod_{\text{odd } i} (1+s_i)$ applied to the regular representation. This is the first layer of the rank hierarchy (cf.\ the author's notes on OEIS A090932 and the descent algebra).

The staircase cascade adds generators one at a time in a specific order ($s_1, s_2, \ldots, s_{n-1}$). The resulting image $\Phi_j(V^H)$ is related to a partial symmetrization that interleaves the $H$-fixed vectors with the new generators.

\section{A Simple Algebraic Proof for the Basic Claim}

\begin{theorem}[Staircase eigenspace property]
  Let $P_j = (1+s_j) \in \mathbb{Q}[S_n]$. Then for any $v \in \mathbb{Q}[S_n]$:
  \[
    s_j \cdot P_j P_{j-1} \cdots P_1 v = P_j P_{j-1} \cdots P_1 v.
  \]
  In particular, $\Phi_j(V^H) \subseteq \ker(s_j - 1)$ for any starting space $V^H$.
\end{theorem}

\begin{proof}
  Let $w = P_{j-1} \cdots P_1 v$ be any element. Then:
  \[
    s_j \cdot P_j w = s_j (1 + s_j) w = (s_j + s_j^2) w = (s_j + 1) w = (1 + s_j) w = P_j w.
  \]
  This holds for any $w$, so $P_j w$ is always in the $+1$ eigenspace of $s_j$.
\end{proof}

\begin{remark}
  The theorem shows the claim holds at the \emph{current} block $j$ trivially: the last factor $(1+s_j)$ creates a $+1$ eigenvector. The more subtle---and more interesting---empirical finding is that \emph{subsequent} blocks do not destroy this property: $\Phi_k(V^H)$ remains in the $+1$ eigenspace of $s_j$ even for $k > j$.
\end{remark}

\section{Stronger Empirical Findings}

\subsection{The cascade does not collapse}

The most surprising computational finding is:

\begin{observation}[Injectivity of cascade]
  For all $n \le 7$ and all $j$:
  \[
    \dim(\Phi_j(V^H)) = \dim(V^H) = \frac{n!}{2^{\lfloor n/2 \rfloor}}.
  \]
\end{observation}

This is not obvious: each $(1+s_j)$ operator has a nontrivial kernel (the $-1$ eigenspace of $s_j$), and one might expect repeated application to gradually collapse the image. That the dimension is preserved throughout the cascade suggests a deep orthogonality between $V^H$ and the kernels of successive $(1+s_j)$ operators.

\subsection{Residual symmetry tail}

Let $S_{\text{tail}}^{(j)} = \{k : s_k \text{ fixes } \Phi_j(V^H)\}$. The computation shows:

\begin{observation}[Residual symmetry tail for $n=6$]
  \[
    S_{\text{tail}}^{(0)} = \{1, 3, 5\}, \quad S_{\text{tail}}^{(j)} = \{j\} \cup (\text{some } k > j) \quad \text{for } j \ge 1.
  \]
  Specifically (for $n=6$):
  \begin{align*}
    S_{\text{tail}}^{(1)} &= \{1, 3, 5\} \\
    S_{\text{tail}}^{(2)} &= \{2, 5\} \\
    S_{\text{tail}}^{(3)} &= \{3, 5\} \\
    S_{\text{tail}}^{(4)} &= \{4\} \\
    S_{\text{tail}}^{(5)} &= \{5\}.
  \end{align*}
\end{observation}

\noindent The generator $s_{n-1} = s_5$ (for $n=6$) persists until step $j=4$ and then vanishes. The pattern suggests that at each step $j$, the ``high-end'' tail generators that commute with all of $s_1, \ldots, s_{j-1}$ (except through higher braid relations) retain their fixing property.

\section{Consequences and Questions}

\subsection{Confirmed claims}
\begin{enumerate}
  \item The staircase eigenspace property (Claim) holds for $n \le 7$, for \emph{all} $j$ (not just odd $j$).
  \item $\dim(\Phi_j(V^H)) = \dim(V^H)$ throughout---the cascade is injective.
  \item The image $\Phi_1(V^H) = V^H$ (block 1 preserves $V^H$), but subsequent blocks move the image.
\end{enumerate}

\subsection{Open questions}
\begin{enumerate}
  \item \textbf{Proof of injectivity:} Why does $\ker(P_j) \cap \Phi_{j-1}(V^H) = \{0\}$ for all $j$? This is the key mystery.
  \item \textbf{Persistence of $s_j$ eigenspace under subsequent blocks:} The computation shows $\Phi_k(V^H) \subseteq \ker(s_j - 1)$ for $k > j$ as well (in many cases). What is the clean characterization?
  \item \textbf{Residual symmetry tail:} What is the precise rule for $S_{\text{tail}}^{(j)}$? The pattern for $n=6, 7$ suggests $s_{n-1}$ fixes the image for all $j \le n-2$, but the pattern for smaller generators is more complex.
  \item \textbf{Irrep-level statement:} Does the claim decompose nicely per-irrep? Does $V_\lambda^H$ (the $H$-fixed part of irrep $V_\lambda$) satisfy the same cascade property restricted to each irrep?
\end{enumerate}

\section{Computational Details}

All computations were performed in Python 3 using \texttt{fractions.Fraction} for exact arithmetic. The regular representation of $S_n$ was used, with permutations stored as tuples and subspaces as lists of sparse dictionaries. Gaussian elimination was performed exactly over $\mathbb{Q}$.

\begin{itemize}
  \item $n=4$: $24 \times 24$ representation, completed in $< 0.01$s.
  \item $n=5$: $120 \times 120$, completed in $0.02$s.
  \item $n=6$: $720 \times 720$, completed in $0.33$s.
  \item $n=7$: $5040 \times 5040$ (sparse), completed in $18.71$s.
\end{itemize}

Script location: \texttt{/home/clio/projects/proofs/staircase\_test\_final.py}

\end{document}
